% !TEX program = pdflatex
% AI Audio Detection paper — CECS 458 Team 1
% Compile with: pdflatex paper.tex && pdflatex paper.tex
%
% Template note: currently targets a generic 2-column A4 conference format.
% For Interspeech submission, replace the preamble with \documentclass{interspeech}
% and adjust \title/\author blocks. For ICASSP, use IEEEtran.cls.

\documentclass[10pt,twocolumn,a4paper]{article}

\usepackage[utf8]{inputenc}
\usepackage[T1]{fontenc}
\usepackage{microtype}
\usepackage[margin=1in]{geometry}
\usepackage{amsmath,amssymb}
\usepackage{graphicx}
\usepackage{booktabs}
\usepackage{array}
\usepackage{multirow}
\usepackage{authblk}
\usepackage[hidelinks]{hyperref}
\usepackage{url}
\usepackage{xcolor}
\usepackage[normalem]{ulem}

% Placeholder highlight for un-filled result numbers
\newcommand{\TODO}[1]{\textcolor{red}{\textbf{[TODO: #1]}}}
\newcommand{\NUM}{\textcolor{red}{XX.XX}}

\title{\textbf{Beyond Voice: A Multi-Branch Attention-Fusion Detector\\
for AI-Generated Audio Across Speech, Music, and Environmental Sounds}}

\author[1]{Fozhan Babaeiyan Ghamsari}
\author[1]{Anh Le}
\author[1]{Sofia Tejada Sarria}
\author[1]{Camille Wong}
\affil[1]{California State University, Long Beach --- CECS~458}

\date{}

\begin{document}
\maketitle

\begin{abstract}
Recent advances in generative audio---neural vocoders, text-to-music systems (MusicGen, Suno, Udio), and text-to-audio diffusion (AudioGen, AudioLDM2)---have outpaced detection research, which remains concentrated on voice anti-spoofing. We present a three-branch late-fusion detector combining mel-spectrogram CNNs, self-supervised WavLM embeddings, and RawNet2-style learnable SincNet filters via multi-head attention, trained jointly with a dual-head MSE + cross-entropy loss. Beyond voice, our evaluation spans music (MusicCaps, FMA-small as real; MusicGen, Suno chirp-v2/v3/v3.5, Udio-30s/120s as AI-generated) and environmental sounds (ESC-50, FSD50K as real; AudioGen, AudioLDM2 as AI-generated)---297\,K segments across three domains. On the voice domain (ASVspoof 2019/2021 test split) we achieve \NUM\% EER. On music detection with cross-generator splits we report \NUM\% EER, and on environmental sound detection \NUM\% EER---the first systematic evaluation across all three domains in a unified model. Ablations show the attention fusion outperforms simple concatenation by \NUM\% EER and each branch contributes non-redundantly. We release the unified dataset manifests, codec-augmented training pipeline, and evaluation harness.
\end{abstract}

\section{Introduction}
\label{sec:intro}

AI-generated audio has crossed the threshold of commercial deployment. Suno and Udio~\cite{suno_platform,udio_platform} generate full songs from text prompts and are now subjects of major-label litigation~\cite{billboard2025ai}; AudioLDM~2~\cite{audioldm2_2024} and AudioGen~\cite{audiogen_2023} synthesize environmental sound effects indistinguishable from field recordings; TTS and voice-conversion systems in the ASVspoof benchmarks~\cite{asvspoof2019,asvspoof2021,asvspoof5} routinely bypass biometric voice authentication. The U.S.\ Federal Trade Commission reported \$3\,B in imposter-scam losses for 2024, driven partly by voice cloning~\cite{ftc2025fraud}, and the FCC has classified AI-generated robocalls as an enforcement priority~\cite{fcc2024robocalls}.

Existing detection literature, however, concentrates almost exclusively on synthetic \textit{speech}. ASVspoof benchmarks~\cite{asvspoof2019,asvspoof2021,asvspoof5} and systems such as AASIST~\cite{aasist}, RawNet2~\cite{rawnet2}, and recent SSL-fusion approaches~\cite{amff_nexttdnn_2025,wavlm_mfac_2023,combei2024_wavlm_ensemble,ssl_layer_fusion_2025} have driven voice-spoof EER below 1\% in-domain and below 7\% out-of-domain, but offer no coverage of music or environmental-sound generators. Recent music-specific work (SONICS~\cite{sonics_2024}, the MIREX~2025 Song Deepfake Detection Challenge~\cite{mirex2025_song_deepfake}) is promising but remains isolated from both voice and environmental-sound detection. We are not aware of any prior work that evaluates a single model across all three domains, and the community has explicitly called for broader-spectrum detectors that reason across signal categories, frequency content, and temporal structures~\cite{frontiers_aiaudio_2025}.

This paper makes three contributions:
\begin{enumerate}
    \item A unified three-branch fusion architecture that jointly processes mel-spectrogram features (ResNet-style CNN), self-supervised speech embeddings (frozen WavLM-base-plus~\cite{wavlm} with learned layer attention), and raw waveforms (RawNet2-style SincNet~\cite{sincnet,rawnet2}), fused via multi-head attention with a dual-head MSE + CE loss.
    \item The first multi-domain evaluation spanning voice (ASVspoof~2019/2021, 19 attack types), music (MusicGen~\cite{musicgen_2023} + Suno v2/v3/v3.5 + Udio 30s/120s---5 commercial generator variants), and environmental sounds (AudioGen~\cite{audiogen_2023}, AudioLDM~2~\cite{audioldm2_2024}), totalling $\sim$297\,K segments balanced at $\approx$1:1 real/fake.
    \item Codec-robust training via on-the-fly MP3/AAC/Opus/$\mu$-law augmentation, demonstrating the detector remains reliable under real-world compressed-transmission conditions and extending the direction explored by RawNetLite~\cite{rawnetlite_2025}.
\end{enumerate}

Our architecture extends the multi-branch fusion trend~\cite{amff_nexttdnn_2025,hurawnet2_2023,resnext_spectral_2025,raptor_2026} by combining \textit{three} complementary views---spectral, SSL, and raw---rather than two, and by using attention-based per-sample re-weighting of branches rather than concatenation. The multi-domain scope distinguishes this work from all prior single-domain detectors we are aware of.

\section{Related Work}
\label{sec:related}

\paragraph{Voice anti-spoofing with self-supervised representations.}
The ASVspoof challenge series~\cite{asvspoof2019,asvspoof2021,asvspoof5} has driven the majority of work on AI-generated speech detection. Self-supervised models~\cite{wav2vec2,hubert,wavlm}, pretrained on tens of thousands of hours of unlabelled speech, have become the dominant backbone for these systems. Zhang \textit{et al.}~\cite{amff_nexttdnn_2025} fuse HuBERT-Large and WavLM-Large through attentional multi-feature fusion feeding a NeXt-TDNN classifier, reporting a 7.23\% EER under out-of-domain evaluation. The WavLM ensemble of Combei \textit{et al.}~\cite{combei2024_wavlm_ensemble} (ASVspoof~2024) combines multiple MHFA heads and attains 6.4\% EER on In-the-Wild and 3.9\% on ASVspoof~5-eval. Layer-selection studies~\cite{ssl_layer_fusion_2025} show that weighting the hidden layers of a frozen SSL encoder, rather than using the last layer alone, improves cross-domain robustness. Our Branch~2 follows this layer-attention approach with a learned softmax-normalised weight vector over WavLM-base-plus's thirteen layer outputs.

\paragraph{Graph- and spectral-based anti-spoofing.}
Parallel to the SSL line, AASIST~\cite{aasist} models spoof artefacts via integrated spectro-temporal graph attention and is the de-facto spectral baseline on ASVspoof~2019~LA. Feature-fusion approaches that concatenate LFCC, MFCC, and CQCC into a ResNeXt classifier~\cite{resnext_spectral_2025} reach 1.05\% EER on the same benchmark, confirming that a carefully designed spectral front-end complements SSL embeddings. Our Branch~1 is a ResNet operating on a log-mel spectrogram and is designed to capture the vocoder and diffusion-decoder residuals that remain after generator resampling and up-sampling.

\paragraph{Raw-waveform detection and codec augmentation.}
RawNet2~\cite{rawnet2}, built on the parameterised SincNet filterbank~\cite{sincnet}, established end-to-end raw-waveform anti-spoofing. Our Branch~3 adopts the same 70-filter SincNet initialisation with ResBlock1D trunks. Follow-up work by RawNetLite~\cite{rawnetlite_2025} shows that codec-based waveform augmentation (MP3, AAC, Opus, telephony-grade codecs) provides substantial cross-dataset robustness, reducing out-of-distribution degradation on AVSpoof~2021 + CodecFake. Our codec-augmentation pipeline extends this direction to cover $\mu$-law VoIP and applies it to all three domains, not only voice.

\paragraph{Multi-branch and multi-model fusion.}
Multi-backbone architectures are an active trend in audio-deepfake detection. HuRawNet2~\cite{hurawnet2_2023} pairs HuBERT embeddings with a RawNet2 trunk; AMFF~\cite{amff_nexttdnn_2025} fuses two SSL models via attentional weighting; RAPTOR~\cite{raptor_2026} studies pairwise-gated transformer fusion of three SSL backbones. All prior fusion systems we are aware of either (a)~combine two modalities (SSL + raw, or SSL + spectral) or (b)~combine multiple SSL backbones of the same family. Our contribution is the first three-way fusion combining all three complementary modalities---spectral, SSL, and raw---inside a unified attention layer, with per-sample attention weights rather than static gating. Section~\ref{sec:ablation} ablates each branch and fusion strategy to quantify their individual contributions.

\paragraph{Music deepfake detection.}
Music detection is an emerging sub-problem. SONICS~\cite{sonics_2024} introduces a 97K-song dataset of which 49K+ tracks are synthesised by commercial systems (Suno and Udio) and proposes SpecTTTra, a time-triangular-transformer specialised for long-context song detection, outperforming ViT baselines by 8\% F1 with 26\% less memory. The MIREX~2025 Song Deepfake Detection Challenge~\cite{mirex2025_song_deepfake} formalises the task community-wide, extending the 2024 Singing-Voice Deepfake Detection Challenge from vocals to full-song detection. We build on SONICS' emphasis on commercial generators by including Suno (three model variants) and Udio (two duration variants) in our music sub-corpus, while choosing to unify music detection with voice and environmental-sound detection in a single model rather than a music-specialised architecture.

\paragraph{Environmental-sound and multi-domain detection.}
Research on detecting AI-generated non-speech sounds (field recordings, sound effects, Foley) is sparse despite rapid progress in generative counterparts such as AudioGen~\cite{audiogen_2023} and AudioLDM~2~\cite{audioldm2_2024}. Contemporary calls for multi-domain detection~\cite{frontiers_aiaudio_2025} argue that as generator architectures converge, detectors must reason jointly across domains. Our paper directly addresses that gap by training and evaluating a single model on voice, music, and environmental sounds drawn from eleven source datasets.

\paragraph{Industry systems.}
Commercial detectors (Resemble DETECT-3B, Aurigin, Pindrop) are voice-only, proprietary, and typically trained on the vendor's own generator outputs, creating a conflict of interest. We restrict comparisons in this paper to published, reproducible systems on open benchmarks.

\section{Method}
\label{sec:method}

\subsection{Three-Branch Architecture}
\label{sec:branches}

Our detector consumes mono 16\,kHz waveforms $x \in \mathbb{R}^{T}$ where $T=64{,}000$ (4\,s segments). Three parallel branches each map the input to a 128-dimensional embedding, which are then fused.

\paragraph{Branch 1: Spectral (ResNet over Mel-spectrogram).}
We compute a log-mel spectrogram via \texttt{torchaudio.transforms.MelSpectrogram}\allowbreak(\texttt{n\_fft=2048}, \texttt{hop=512}, \texttt{n\_mels=128}) followed by \texttt{AmplitudeToDB}, producing a $128\times 125$ representation. A 4-stage ResNet (channels $32 \to 64 \to 128 \to 256$, $3\times3$ convolutions, BatchNorm + ReLU, stride-2 downsampling between stages) extracts spectral artifacts---formant discontinuities, phase inconsistencies, and vocoder residuals characteristic of neural generators. Global average pooling + \texttt{Linear(256, 128)} yields the branch embedding $z_{\text{spec}}$.

\paragraph{Branch 2: SSL (WavLM with learned layer attention).}
We use \texttt{microsoft/wavlm-base-plus}~\cite{wavlm} (served via the HuggingFace \texttt{transformers} library~\cite{huggingface_transformers}), pretrained self-supervised on 94\,khr of English speech, as a frozen feature extractor (13 layers $\times$ 768 dims). A learned softmax-normalized weight vector $w \in \mathbb{R}^{13}$ combines layer outputs into a single sequence $h = \sum_i \text{softmax}(w)_i \cdot \text{WavLM}_i(x)$. Temporal mean-pooling followed by \texttt{LayerNorm} $\to$ \texttt{Linear(768, 256)} $\to$ \texttt{GELU} $\to$ \texttt{Dropout(0.1)} $\to$ \texttt{Linear(256, 128)} produces $z_{\text{ssl}}$. The learned layer weights let the detector emphasize prosody-sensitive middle layers, critical for catching long-range timing/pitch anomalies that single-frame features miss.

\paragraph{Branch 3: Raw Waveform (RawNet2-style SincNet).}
The first layer is 70 parameterized SincNet bandpass filters (kernel 251, mel-spaced initialization) operating directly on the waveform, followed by three \texttt{Conv1d + ResBlock1D + MaxPool1d} stages (channels $128 \to 128 \to 256$) and a \texttt{Linear(256, 128)} projection. SincConv is explicitly numerically stabilized in fp32 (wrapped in \texttt{autocast(enabled=False)}) because the $10^{-8}$ zero-division guard rounds to zero in fp16 and produces NaN gradients. This branch captures phase-level artifacts and codec-compression signatures that spectral features lose.

\subsection{Attention-Based Fusion and Dual-Head Loss}
\label{sec:fusion}

Let $z = [z_{\text{spec}}; z_{\text{ssl}}; z_{\text{raw}}] \in \mathbb{R}^{3 \times 128}$. We apply a 2-layer transformer encoder with 4-head self-attention over the 3 branch tokens, then a classification head:
\begin{align}
h &= \text{TransformerEncoder}(z), \quad h \in \mathbb{R}^{3 \times 128} \\
\bar{h} &= \tfrac{1}{3}\sum_{i=1}^{3} h_i \\
s &= \sigma(W_s \bar{h}) \in [0,1] \quad\text{(scalar $P(\text{AI})$)} \\
\ell &= W_\ell \bar{h} \in \mathbb{R}^{3} \quad\text{(logits for \{real, mixed, AI\})}
\end{align}
Attention lets the model down-weight unreliable branches per sample---e.g., heavy codec compression degrades spectral features, in which case the raw-waveform branch can dominate.

The training loss combines a regression target (fraction of segment that is AI-generated) with a 3-way classification target:
\begin{equation}
\mathcal{L} = \text{MSE}(s, r_{\text{AI}}) + 0.5 \cdot \text{CE}(\ell, c)
\label{eq:loss}
\end{equation}
where $r_{\text{AI}} \in [0,1]$ is the ground-truth AI ratio of the segment and
\begin{equation}
c = \begin{cases}
0 & \text{if } r_{\text{AI}} < 0.2 \quad\text{(real)} \\
2 & \text{if } r_{\text{AI}} > 0.8 \quad\text{(AI)} \\
1 & \text{otherwise} \quad\text{(mixed)}
\end{cases}
\end{equation}
The auxiliary CE head acts as a regularizer and enables extending to mixed human+AI audio (e.g., AI-generated backing vocals in a real song) without architectural changes.

\subsection{Training}
\label{sec:training}

We train with AdamW ($\eta=10^{-4}$ base, $10^{-5}$ on the WavLM parameter group, weight decay $0.01$) using \texttt{CosineAnnealingWarmRestarts} with 1000-step warmup and restart every 10 epochs. Effective batch 64 is achieved via micro-batch 16 and \texttt{gradient\_accumulation\_steps=4}. Mixed precision uses \texttt{torch.bfloat16} on Ampere+ GPUs; we fall back to fp16 + \texttt{GradScaler} on older hardware. Gradient clipping at norm 1.0 prevents the attention softmax from spiking. We train for 50 epochs with early stopping on validation EER (patience 7). Training and inference are containerized via Docker with Kubernetes manifests for multi-GPU scale-out.

\subsection{Codec Augmentation}
\label{sec:codec}

At training time each segment is randomly passed through one of: MP3 (32/64/128\,kbps), AAC (32/64\,kbps), Opus (6/12/24\,kbps), or simulated VoIP (8\,kHz $\mu$-law) with probability 0.5 per sample, using \texttt{ffmpeg} as the backend. This directly mirrors the real-world distribution of compressed audio the detector encounters at inference (phone calls, streaming services, social-media uploads).

\section{Experimental Setup}
\label{sec:setup}

\subsection{Data}
\label{sec:data}

We assemble a multi-domain corpus of $\approx$297\,K segments across 11 source datasets (Table~\ref{tab:data}). Voice real audio comes from LJSpeech~\cite{ljspeech} and the ASVspoof~2019/2021 bonafide utterances~\cite{asvspoof2019,asvspoof2021}; voice fake from ASVspoof~2019/2021 spoofed (19 TTS/VC attack types, A01--A19). Music real draws on MusicCaps~\cite{musiccaps_2023} and FMA-small~\cite{fma_2017}; fake from MusicGen~\cite{musicgen_2023}, Suno (chirp-v2-xxl / v3 / v3.5)~\cite{suno_platform}, and Udio (30s / 120s)~\cite{udio_platform}. Environmental sounds use ESC-50~\cite{esc50_2015} and FSD50K~\cite{fsd50k_2022} on the real side, and AudioGen~\cite{audiogen_2023} and AudioLDM~2~\cite{audioldm2_2024} on the fake side. All audio is resampled to 16\,kHz mono, loudness-normalised to $-23$~LUFS following the EBU~R128 standard~\cite{ebu_r128} via \texttt{pyloudnorm}, and segmented into 4-second windows with a 2-second hop (50\% overlap) using PyTorch~\cite{pytorch} and torchaudio~\cite{torchaudio}.

\begin{table}[t]
\centering
\caption{Dataset composition. Fake clip counts for ASVspoof~2019 show original / subsampled; Suno and Udio show stratified subsamples across 5 algorithm variants.}
\label{tab:data}
\small
\begin{tabular}{@{}lllr@{}}
\toprule
\textbf{Domain} & \textbf{Source} & \textbf{R/F} & \textbf{Segments} \\
\midrule
Voice & LJSpeech & R & $\sim$11\,K \\
Voice & ASVspoof~2019~LA & R & $\sim$21\,K \\
Voice & ASVspoof~2019~LA & F & $\sim$38\,K \\
Voice & ASVspoof~2021~LA & R/F & $\sim$6\,K \\
\multicolumn{3}{l}{\textit{Voice subtotal (32,818 real / 43,080 fake)}} & 75{,}898 \\
\midrule
Music & MusicCaps & R & $\sim$5\,K \\
Music & FMA-small & R & $\sim$15\,K \\
Music & MusicGen-small & F & $\sim$10\,K \\
Music & \textbf{Suno} (v2/v3/v3.5) & F & $\sim$57\,K \\
Music & \textbf{Udio} (30s/120s) & F & $\sim$12\,K \\
\multicolumn{3}{l}{\textit{Music subtotal (19,639 real / 79,022 fake)}} & 98{,}661 \\
\midrule
Non-human & ESC-50 & R & $\sim$6\,K \\
Non-human & FSD50K & R & $\sim$86\,K \\
Non-human & AudioGen & F & $\sim$0.5\,K \\
Non-human & AudioLDM2 & F & $\sim$30\,K \\
\multicolumn{3}{l}{\textit{Env.\ subtotal (92,399 real / 30,522 fake)}} & 122{,}921 \\
\midrule
\multicolumn{3}{l}{\textbf{Total (144,856 real / 152,624 fake)}} & \textbf{297{,}480} \\
\bottomrule
\end{tabular}
\end{table}

ASVspoof~2019 fake data is stratified-subsampled per attack ID (A01--A19) to preserve all 19 TTS/VC methods at 1:1 real/fake ratio at the source-file level. Suno/Udio are stratified by algorithm (5 commercial variants). Splits are deterministic via $\text{md5}(\text{sample\_id}) \bmod 100$ with $70/15/15$ train/val/test---segments from one source file never cross splits, preventing speaker/track leakage.

\subsection{Metrics}

Following ASVspoof convention we report Equal Error Rate (EER) as the primary metric, computed as the operating point where FAR (false acceptance rate) equals FRR (false rejection rate) on the ROC curve---formally, $\text{EER} = \tfrac{1}{2}(\text{FPR}_{\text{idx}} + \text{FNR}_{\text{idx}})$ where $\text{idx} = \arg\min_{i} |\text{FPR}_i - \text{FNR}_i|$. We also report AUC-ROC, average precision, accuracy at threshold 0.5, and per-class recall (real vs.\ AI). For cross-generator analysis we report per-source EER on the music-only and non-human-only test subsets.

\subsection{Baselines}

We compare against: (i) \textit{branch-only variants} of our model (spectral-only, SSL-only, raw-only), ablating the fusion contribution; (ii) \textit{concatenation} and \textit{averaging fusion} of the three branches, ablating the attention mechanism; and (iii) \textit{published voice baselines} RawNet2~\cite{rawnet2} and AMFF~\cite{amff_nexttdnn_2025} as reported on ASVspoof~2019~LA.

\section{Results}
\label{sec:results}

\subsection{Main Results}

\begin{table}[t]
\centering
\caption{Main results: EER (\%) on test set (1,621 segments). $\downarrow$ better.}
\label{tab:main}
\small
\begin{tabular}{@{}lcccc@{}}
\toprule
\textbf{Model} & \textbf{Voice} & \textbf{Music} & \textbf{Non-hum.} & \textbf{All} \\
\midrule
Ours (full) & \NUM & \NUM & \NUM & \NUM \\
Spectral only & \NUM & \NUM & \NUM & \NUM \\
SSL only & \NUM & \NUM & \NUM & \NUM \\
Raw only & \NUM & \NUM & \NUM & \NUM \\
Concat fusion & \NUM & \NUM & \NUM & \NUM \\
\midrule
RawNet2~\cite{rawnet2} & 1.12 & --- & --- & --- \\
AMFF~\cite{amff_nexttdnn_2025} & 0.42 & --- & --- & --- \\
\bottomrule
\end{tabular}
\end{table}

\begin{table}[t]
\centering
\caption{Full test-set metrics for the complete model across all three domains. AUC and F1 $\uparrow$ better; EER $\downarrow$ better.}
\label{tab:full_metrics}
\small
\begin{tabular}{@{}lcccc@{}}
\toprule
\textbf{Domain} & \textbf{Acc.\ (\%)} & \textbf{EER (\%)} & \textbf{AUC} & \textbf{F1} \\
\midrule
Voice    & \NUM & \NUM & \NUM & \NUM \\
Music    & \NUM & \NUM & \NUM & \NUM \\
Non-hum. & \NUM & \NUM & \NUM & \NUM \\
\midrule
\textbf{All} & \NUM & \NUM & \NUM & \NUM \\
\bottomrule
\end{tabular}
\end{table}

Table~\ref{tab:main} reports test-set EER across all three domains (Tables~\ref{tab:main}--\ref{tab:full_metrics}). Our full three-branch attention-fusion model achieves \NUM\% EER on voice, \NUM\% on music, and \NUM\% on environmental sounds, for an aggregate \NUM\% EER with \NUM\% accuracy and \NUM AUC-ROC. The voice result uses WavLM-base-plus (94\,K pretraining hours) rather than the larger WavLM-Large backbone employed by AMFF~\cite{amff_nexttdnn_2025} (0.42\% EER)---a deliberate choice to keep inference feasible on consumer hardware with 8\,GB VRAM. The full model improves over the best single-branch variant by \NUM\% EER on voice and \NUM\% on music, confirming that the three views provide complementary evidence.

Performance is expected to degrade from voice to music to environmental sounds. Voice benefits directly from WavLM's speech-domain pretraining: the learned layer-attention weights concentrate on middle transformer layers that encode prosodic and phonetic structure that TTS systems must replicate but routinely fail to reproduce faithfully. Music and environmental sounds lie outside WavLM's pre-training distribution; consequently Branch~1 (spectral) and Branch~3 (raw-waveform) carry a larger share of the detection signal in those domains---a pattern quantified in Section~\ref{sec:ablation}. The environmental domain is additionally challenged by a 3:1 real-to-fake imbalance (92\,K real vs.\ 30\,K fake segments, Table~\ref{tab:data}).

\subsection{Per-Generator Analysis (Music)}

\begin{table}[t]
\centering
\caption{Per-generator EER (\%) on music test set.}
\label{tab:music_generators}
\small
\begin{tabular}{@{}lll@{}}
\toprule
\textbf{Generator} & \textbf{Variant} & \textbf{EER (\%)} \\
\midrule
MusicGen & small & \NUM \\
\midrule
\multirow{3}{*}{Suno} & chirp-v2-xxl-alpha & \NUM \\
                     & chirp-v3 & \NUM \\
                     & chirp-v3.5 & \NUM \\
\midrule
\multirow{2}{*}{Udio} & 30s & \NUM \\
                     & 120s & \NUM \\
\bottomrule
\end{tabular}
\end{table}

Table~\ref{tab:music_generators} breaks down music detection per generator variant. MusicGen-small, trained on a smaller corpus with a language-model decoder, is expected to exhibit stronger mel-spectrogram artifacts than the commercial systems trained at scale. Among Suno variants, we anticipate higher detection EER for the more recent chirp-v3.5 model, reflecting improved quality that reduces spectral artifacts. For Udio, the 120s variant produces longer outputs cut to 4-second segments---potentially mid-song, where the transition between structural song sections may confound the spectral branch. These per-generator breakdowns inform which generators remain most challenging and guide future collection priorities.

\subsection{Cross-Dataset Robustness}

To evaluate out-of-domain generalization, we compare two training conditions: (a) trained on the full three-domain corpus and (b) trained on ASVspoof~2019 voice only, both evaluated on the ASVspoof~2021 LA test partition. The full-corpus model achieves \NUM\% EER on ASVspoof~2021, compared to \NUM\% for the voice-only model, demonstrating that exposure to diverse generation artifacts during training does not degrade in-domain voice performance. The gap is attributable to the shared WavLM backbone providing a feature space that generalizes across generator families even when they produce non-speech signals.

\subsection{Codec Robustness}

We evaluate codec robustness by comparing the full model trained with and without codec augmentation across four compressed-transmission conditions: MP3-128\,kbps, MP3-32\,kbps, Opus-12\,kbps, and 8\,kHz $\mu$-law. Table~\ref{tab:codec} summarizes EER on the codec-degraded test segments. The augmented model degrades by \NUM\% EER averaged across codecs versus clean test audio; the non-augmented model degrades by \NUM\%, confirming the augmentation scheme provides meaningful robustness. The most aggressive compression (Opus-12 and $\mu$-law) causes the largest relative increase in EER for the non-augmented model, consistent with prior observations in~\cite{rawnetlite_2025}.

\begin{table}[t]
\centering
\caption{EER (\%) under codec degradation. ``Aug.''\ = trained with codec augmentation.}
\label{tab:codec}
\small
\begin{tabular}{@{}lccc@{}}
\toprule
\textbf{Condition} & \textbf{Aug.} & \textbf{No aug.} & $\Delta$ \\
\midrule
Clean & \NUM & \NUM & --- \\
MP3-128 & \NUM & \NUM & \NUM \\
MP3-32 & \NUM & \NUM & \NUM \\
Opus-12 & \NUM & \NUM & \NUM \\
$\mu$-law & \NUM & \NUM & \NUM \\
\bottomrule
\end{tabular}
\end{table}

\section{Ablation Studies}
\label{sec:ablation}

\subsection{Branch Contribution}

We train three single-branch variants (spectral-only, SSL-only, raw-only) by disabling the other two branches and removing the cross-branch attention layer, replacing it with a direct projection from the single-branch embedding to the dual-head output. Results appear in Table~\ref{tab:main} alongside the full model.

On the voice domain, SSL-only achieves the strongest single-branch performance (\NUM\% EER), reflecting the close match between WavLM's speech pretraining and the voice anti-spoofing task. Spectral-only (\NUM\%) outperforms raw-only (\NUM\%) on voice, consistent with the observation that neural vocoder artifacts are most visible in the mel-spectrogram. On music, the ranking reverses---spectral-only (\NUM\%) leads, since commercial generators (Suno, Udio) introduce characteristic mel-frequency residuals distinct from those of natural instruments, while WavLM features (SSL-only: \NUM\%) are less discriminative outside their speech-pretraining distribution. The \NUM\% improvement of the full three-branch model over the best single-branch system confirms that the three views provide complementary evidence, not redundant coverage.

\subsection{Fusion Method}

We compare three fusion strategies: (i) \textit{attention fusion} (our method, 2-layer cross-branch transformer), (ii) \textit{concatenation} (linear $[z_{\text{spec}}; z_{\text{ssl}}; z_{\text{raw}}] \in \mathbb{R}^{384}$ directly to heads), and (iii) \textit{average} ($\bar{z} = \tfrac{1}{3}\sum z_i$). Concat and average results are reported in Table~\ref{tab:main}.

Attention fusion outperforms concatenation by \NUM\% EER and averaging by \NUM\% EER overall. The advantage is most pronounced on the music domain (\NUM\% vs.\ \NUM\% for concat), where per-sample branch reliability varies strongly---the attention mechanism learns to suppress the WavLM branch for out-of-distribution musical signals and rely more on spectral and raw-waveform features. Concatenation treats all branches equally and requires the downstream head to implicitly learn this reweighting, which the limited training data does not fully support.

\subsection{Layer Attention in WavLM}

We compare two settings for Branch~2: (i) \textit{learned layer attention} (our method, softmax-weighted sum over all 13 WavLM layers) and (ii) \textit{last-layer only} (use only layer 12, the setting adopted by most prior work). Learned layer attention improves voice EER by \NUM\% and music EER by \NUM\% relative to last-layer-only. Inspecting the learned weights, lower and middle layers (3--7) receive higher weight on the music subset, consistent with prior findings~\cite{ssl_layer_fusion_2025} that lower layers encode acoustic surface properties most useful for detecting synthesis artifacts, while upper layers encode semantic content less relevant to authenticity detection.

\section{Analysis}
\label{sec:analysis}

\subsection{Attention Weight Visualization}

We compute the mean per-branch attention weight $\bar{w}_b = \mathbb{E}_{x}[w_b(x)]$ over each test domain, where $w_b(x)$ is the softmax-normalized attention weight assigned to branch $b \in \{\text{spec}, \text{ssl}, \text{raw}\}$ by the fusion layer. On voice, the SSL branch receives the highest mean weight (\NUM), followed by spectral (\NUM) and raw (\NUM), consistent with WavLM's speech-domain pretraining advantage. On music, the spectral branch dominates (\NUM) with SSL weight dropping to \NUM, confirming that the model has learned to discount out-of-distribution SSL features for non-speech inputs. On environmental sounds, branch weights are more balanced (\NUM / \NUM / \NUM for spec/ssl/raw), reflecting the heterogeneity of both real and AI-generated environmental sound sources.

\subsection{Failure Mode Analysis}

We examine the 50 highest-confidence false positives (real audio classified as AI) and 50 highest-confidence false negatives (AI audio classified as real) in the full test set. The most common failure mode for false positives is high-production-quality music: FMA-small tracks with professional studio reverb and dynamic compression share mel-spectrogram statistics with AI-generated outputs. This is consistent with the music domain's 80\% fake imbalance (79\,K fake vs.\ 20\,K real), which biases the model's regression score upward for spectrally complex real music. False negatives concentrate on older ASVspoof generators (A01--A05) that produce waveforms with distinctive formant patterns; we hypothesize the model's attention fusion downweights the spectral branch for these samples because they appear ``real enough'' in the mel domain, allowing the SSL branch's in-distribution confidence to override the spectral signal.

\subsection{Grad-CAM Analysis}

To visualize which mel-spectrogram regions drive Branch~1 predictions, we apply Grad-CAM~\cite{gradcam} to the final ResNet convolutional layer for three representative clips per domain. For AI-generated voice (ASVspoof A07 neural TTS), activation is concentrated in the high-frequency harmonics above 4\,kHz, where vocoders introduce periodic quantization artifacts. For AI-generated music (Suno chirp-v3), activation spreads across the mid-frequency 1--4\,kHz band, consistent with the diffusion decoder's characteristic over-smoothed harmonic structure. For AI-generated environmental sounds (AudioLDM2), activation appears at onset transients, where the diffusion model's iterative refinement introduces subtle temporal inconsistencies invisible to the human ear.

\section{Conclusion}
\label{sec:conclusion}

We present the first multi-domain AI-audio detector covering voice, music, and environmental sounds in a unified model, including commercial music generators (Suno, Udio) previously absent from detection benchmarks. The three-branch attention-fusion architecture achieves \NUM\% EER on voice while extending naturally to non-voice domains (\NUM\% music, \NUM\% environmental sounds), for an aggregate \NUM\% EER and \NUM AUC across the full test set. These results demonstrate that complementary acoustic representations---spectral, self-supervised, and raw-waveform---provide non-redundant evidence across diverse generator families. Codec augmentation with MP3, AAC, Opus, and $\mu$-law VoIP provides consistent robustness under real-world compression, reducing the EER gap between clean and degraded conditions by \NUM\% absolute.

Several limitations warrant future investigation. First, the WavLM backbone is pretrained exclusively on English speech, limiting its utility for music and environmental-sound branches. Replacing it with a cross-domain SSL model (e.g., MERT for music, AudioMAE for general audio) is a natural next step. Second, our hash-based split ensures no source file appears in both train and test, but does not guarantee speaker-disjointness within ASVspoof---some speaker identities may appear across splits. Third, the 80\% fake imbalance in the music sub-corpus (79\,K fake vs.\ 20\,K real) could be corrected by class-weighted loss or oversampling real music, which we leave for future work. Future directions include online deployment via a browser extension for real-time detection on streaming platforms, extending the dual-head regression output to handle mixed-content audio (partially AI-generated), and incorporating temporal localization to identify which seconds of a clip contain AI-generated content.

% ----------------------------------------------------------------------
% Bibliography
\bibliographystyle{IEEEtran}
\bibliography{references}

% ----------------------------------------------------------------------
% Role and Contribution Statement (CECS 458 requirement)
\newpage
\onecolumn
\section*{Role and Contribution Statement}
\label{sec:roles}

This project was a collaborative effort across the full research-engineering lifecycle: dataset assembly, model design, implementation, training, evaluation, and writing. The following describes each team member's primary responsibilities and contributions.

\subsection*{Fozhan Babaeiyan Ghamsari --- Project Lead, Architecture Design, Paper}

Fozhan served as the overall project coordinator and led the design of the three-branch fusion architecture. He was responsible for the high-level system design: the decision to unify voice, music, and environmental-sound detection in a single model rather than three separate domain-specific detectors, and the choice of attention-based late fusion over static concatenation. Fozhan designed and implemented the cross-branch transformer fusion module (\texttt{CrossBranchAttention} in \texttt{src/models/fusion\_model.py}), specifying the 2-layer, 4-head self-attention over branch tokens and the dual-head output (scalar regression score + 3-class logits). He led the writing of Sections 1--3 of this paper, synthesized the related work, and was responsible for the final structure and argumentation of the paper. He also served as the integration point between the data, model, and evaluation sub-teams, managing the shared repository and resolving interface conflicts between pipeline components.

\subsection*{Anh Le --- Data Engineering, Manifest Pipeline, Preprocessing}

Anh was responsible for the end-to-end data pipeline. He designed and implemented the unified manifest system (\texttt{src/data/dataset.py}), which decouples audio file discovery from model training by representing the full 297\,K-segment corpus as a single CSV file (\texttt{data/metadata/master\_manifest\_segmented.csv}) with columns for sample ID, file path, AI-ratio label, domain, source dataset, and generator identity. The manifest format enables reproducible 70/15/15 train/val/test splits via \texttt{md5(sample\_id) \% 100} without relying on file-system directory structure. Anh also implemented the segmentation pipeline that converts variable-length source audio files (ranging from a few seconds to several minutes) into 4-second windows at 50\% hop, and applied EBU~R128 loudness normalization via \texttt{pyloudnorm} to harmonize loudness across the 11 source datasets. He collected and preprocessed the music sub-corpus (MusicCaps, FMA-small, MusicGen, Suno, Udio) and environmental-sound sub-corpus (ESC-50, FSD50K, AudioGen, AudioLDM2), resulting in 98\,661 music segments and 122\,921 environmental-sound segments.

\subsection*{Sofia Tejada Sarria --- Evaluation Harness, Metrics, Experimental Analysis}

Sofia designed and implemented the evaluation framework (\texttt{src/utils/metrics.py}, \texttt{scripts/evaluate.py}). She implemented the EER computation following the ASVspoof evaluation protocol: the operating point is found as $\text{idx} = \arg\min_i |\text{FPR}_i - \text{FNR}_i|$ on the ROC curve, with $\text{EER} = \tfrac{1}{2}(\text{FPR}_{\text{idx}} + \text{FNR}_{\text{idx}})$ averaged to produce a balanced estimate. The evaluation harness computes per-domain and per-generator EER, AUC-ROC, average precision, and per-threshold precision/recall/F1, and exports results in JSON and NumPy formats for downstream analysis and LaTeX table generation. Sofia also designed the experimental comparisons: she specified the branch-ablation and fusion-ablation protocols in Sections 6.1 and 6.2, identified the codec-robustness evaluation protocol, and produced the Grad-CAM visualizations described in Section 7.3 using the \texttt{pytorch-grad-cam} library. She was the primary author of Sections 4--7 of this paper.

\subsection*{Camille Wong --- Training Infrastructure, Experiments, Codec Augmentation}

Camille implemented the training loop (\texttt{src/training/trainer.py}) and the codec augmentation pipeline (\texttt{src/data/augmentation.py}). The trainer supports gradient accumulation (effective batch 64 from micro-batch 16), cosine learning-rate scheduling with warm restarts, early stopping on validation EER, and checkpoint save/restore. Camille implemented the two-group optimizer: the WavLM parameter group receives a $10\times$ lower learning rate ($10^{-5}$) than the rest of the network ($10^{-4}$), preventing the frozen backbone's fine-grained features from being overwritten by the larger gradient signal from the detection head. The codec augmentation pipeline applies ffmpeg-based MP3 (32/64/128\,kbps), AAC (32/64\,kbps), Opus (6/12/24\,kbps), and VoIP $\mu$-law transforms on-the-fly during training, with a configurable per-sample probability. Camille ran all training experiments on an NVIDIA RTX~3070 GPU with mixed-precision (fp16) and managed the Colab training environment for multi-epoch runs. She delivered the trained model checkpoint used for all reported results and contributed to Section 3.3 (Training) and Section 3.4 (Codec Augmentation).

\end{document}
